\begin{abstractpage}


\begin{abstract}{english}
Implementing robust Machine Learning (ML) solutions in medical settings is a challenging task due to a number of factors, such as the difficulty of acquiring high-quality labeled data, the lack of sufficiently large datasets, the complexity of medical data, the need for interpretability and explainability, and the strict regulatory requirements. In this thesis, we aim to address some of these challenges in a specific medical application: the classification of glioma tumors using Raman Spectroscopy (RS) data. We analyze a dataset of 58 tumor biopsies, and adress the challenges of high-dimensionality, small sample size, and class imbalance. We propose a generative data augmentation approach based on a Conditional Variational Autoencoder ($\beta$-CVAE) architecture to generate synthetic RS data, and we evaluate the classification performance of several ML models under various regimes, namely: (i) Train on Real Test on Real (TR/TR), (ii) Train on Synthetic Test on Real (TS/TR), and (iii) Train on Real + Synthetic Test on Real (TSR/TR). We show that the proposed generative data augmentation approach can act as a regularizer and improve the classification performance of ML models. \TODO{Add the headline numbers -- an abstract that reports no findings undersells the work. The figures, from Table 6.1 (binary, \texttt{exp\_201}) and Table 6.2 (six-class, \texttt{exp\_204}), both slide-level balanced accuracy: binary TR/TR peaks at 0.742 (XGBoost); TS/TR drops to 0.622--0.669 for the tree models and further for the others; TSR/TR shifts the baseline up by +0.004 to +0.014, consistently in sign but with every 95\% CI including zero. The six-class task reaches 0.289 against a chance level of 0.167. State the TSR/TR result as a tendency, not a demonstrated effect -- that is what the CIs support. Worth one clause on the mechanism found in Chapter 7: the generator reproduces class-mean spectra faithfully but collapses their variance.}
\end{abstract}

\end{abstractpage}